\documentclass{article}
\usepackage{amsmath, amssymb, amsthm}
\usepackage{hyperref}
\usepackage{geometry}
\geometry{margin=1in}

\title{Erd\H{o}s Problem \#791}
\date{Last updated: 2025-08-31}
\author{Source: erdosproblems.com}

\begin{document}
\maketitle

\noindent\textbf{Status:} OPEN \quad \textbf{No prize}

\noindent\textbf{Tags:} additive combinatorics

\medskip

\noindent\textbf{Problem Statement:}

Let $g(n)$ be minimal such that there exists $A\subseteq \{0,\ldots,n\}$ of size $g(n)$ with $\{0,\ldots,n\}\subseteq A+A$. Estimate $g(n)$. In particular is it true that $g(n)\sim 2n^{1/2}$?




Such a set is often called a finite additive $2$-basis. A problem of Rohrbach, who proved in \cite{Ro37}\[(2+c)n \leq g(n)^2 \leq 4n\]for some small constant $c&#62;0$. The current best-known bounds are\[(2.181\cdots+o(1))n\leq g(n)^2 \leq (3.458\cdots+o(1))n.\]The lower bound is due to Yu \cite{Yu15}, and the upper bound is due to Kohonen \cite{Ko17}. (The disproof of $g(n)\sim 2n^{1/2}$ was accomplished by Mrose \cite{Mr79}, who gave a construction implying $g(n)^2 \leq \frac{7}{2}n$.)




References

[Ko17] Kohonen, Jukka, An improved lower bound for finite additive 2-bases . J. Number Theory (2017), 518--524.

[Mr79] Mrose, Arnulf, Untere Schranken f\"ur die {R}eichweiten von {E}xtremalbasen fester {O}rdnung . Abh. Math. Sem. Univ. Hamburg (1979), 118--124.

[Ro37] Rohrbach, Hans, Ein {B}eitrag zur additiven {Z}ahlentheorie . Math. Z. (1937), 1--30.

[Yu15] Yu, Gang, A new upper bound for finite additive {$h$}-bases . J. Number Theory (2015), 95--104.

\noindent\textbf{Related OEIS sequences:} A066063


\bigskip
\noindent\small{Source: \url{https://www.erdosproblems.com/791}}
\end{document}
